%% bab1.tex
%% BAB 1 - PENDAHULUAN
%% Format sesuai Panduan Penyusunan Tesis ITS 2021

\chapter{PENDAHULUAN}

\section{Latar Belakang}

Latar belakang menyajikan konteks penelitian, untuk apa penelitian ini dilakukan, dan hal apa yang mengarahkan penelitian ini. Di sini diuraikan dalam keadaan bagaimana topik akan dilakukan. Latar belakang memuat studi awal atau berbagai teori utama yang relevan dan baru (\textit{recent}) yang dipadukan sehingga mengerucut pada suatu persoalan unik yang kemudian disusun dalam bentuk perumusan masalah.

Lazimnya bagian ini diawali dengan menguraikan kesenjangan, baik secara teoretis maupun praktis, antara harapan dan kenyataan. Setiap paragraf dimulai dengan masuk satu tab (1,5 cm) dari batas margin kiri. Antar paragraf tidak perlu diberi jarak tambahan.

Untuk mengutip referensi, gunakan sistem Harvard yaitu dengan mencantumkan nama belakang penulis dan tahun penerbitan dalam tanda kurung, misalnya \citep{contoh2024}. Apabila penulis dari artikel yang dikutip lebih dari dua orang maka cukup nama penulis pertama yang ditulis, kemudian dilanjutkan dengan tulisan dkk.

\section{Perumusan Masalah}

Dalam sub-bab ini, permasalahan yang ingin diselesaikan dirumuskan secara jelas, tajam, dan terfokus. Bagian ini memuat uraian/pernyataan atau berbagai topik pokok yang akan digali dalam penelitian ini. Definisi, asumsi, dan lingkup penelitian/studi dapat pula dijelaskan pada bagian ini.

\begin{enumerate}
    \item Pertanyaan penelitian pertama?
    \item Pertanyaan penelitian kedua?
    \item Pertanyaan penelitian ketiga?
\end{enumerate}

\section{Tujuan Penelitian}

Pada bagian ini, tujuan dilakukannya penelitian/studi dan target atau sasaran yang ingin dicapai dinyatakan secara singkat dan jelas sesuai dengan permasalahan yang telah dirumuskan. Tujuan penelitian ini adalah:

\begin{enumerate}
    \item Tujuan pertama sesuai dengan rumusan masalah pertama.
    \item Tujuan kedua sesuai dengan rumusan masalah kedua.
    \item Tujuan ketiga sesuai dengan rumusan masalah ketiga.
\end{enumerate}

\section{Manfaat Penelitian}

Di sini perlu dicantumkan manfaat/kegunaan khusus/dampak kemanfaatan yang diharapkan dari hasil penelitian/studi ini. Manfaat penelitian dapat dibagi menjadi manfaat teoretis dan manfaat praktis.

\subsection{Manfaat Teoretis}

Secara teoretis, penelitian ini diharapkan dapat memberikan kontribusi berupa:

\begin{enumerate}
    \item Kontribusi teoretis pertama.
    \item Kontribusi teoretis kedua.
\end{enumerate}

\subsection{Manfaat Praktis}

Secara praktis, penelitian ini diharapkan dapat memberikan manfaat berupa:

\begin{enumerate}
    \item Manfaat praktis pertama.
    \item Manfaat praktis kedua.
\end{enumerate}

\section{Batasan Masalah}

Batasan masalah berisi semua variabel yang diteliti atau kondisi yang melingkupi penelitian. Dengan menampilkan bagian ini, pembaca dapat menyikapi dengan tepat laporan penelitian yang disajikan. Batasan masalah dalam penelitian ini adalah:

\begin{enumerate}
    \item Batasan pertama.
    \item Batasan kedua.
    \item Batasan ketiga.
\end{enumerate}
